\chapter*{Preface: How This Selection Was Made}

This booklet distills a working research library of \textbf{99 lattice-QCD
articles} (the \texttt{papers/} collection) into the set of papers that
matter most. The selection draws on two and only two sources, as required:

\begin{enumerate}
  \item the \textbf{original articles themselves}, several of which are
        classics of the field in their own right; and
  \item their \textbf{reference lists}: every bibliography of every article
        was mined, and the citation counts were tallied across the whole
        corpus.
\end{enumerate}

A paper earns its place here for one (or both) of two reasons:
(i)~it is an \emph{article of the library itself} whose content is a
milestone of the field, or (ii)~it is a \emph{hub of the citation network}
--- a foundational work that the library articles lean on again and again.
The citation statistics quoted below are conservative estimates from this
particular corpus; they measure relevance \emph{to this library}, not
global impact (which would only be larger).

\section*{The map of the territory}

The 38 selected papers organize themselves naturally into eight chapters,
which are also the eight load-bearing pillars of modern lattice parton
physics:

\begin{center}
\begin{tabular}{clc}
\toprule
Ch. & Theme & Papers\\
\midrule
1 & QCD and the parton idea: theoretical foundations & 3\\
2 & Simulation engines and signal enhancement & 4\\
3 & The Yang--Mills gradient flow program & 6\\
4 & From Euclidean lattices to light cones: LaMET & 7\\
5 & The renormalization battle & 7\\
6 & Gluon PDFs: where the two lines converge & 5\\
7 & TMDs and global fitting: contact with phenomenology & 3\\
8 & Stochastic quantization reborn: AI sampling & 3\\
\bottomrule
\end{tabular}
\end{center}

Two storylines run through the whole book and meet in Chapter~6:
\begin{itemize}
  \item the \textbf{gradient-flow line} (L\"uscher's program): smooth
        gauge fields as renormalized fields, from trivializing maps to
        smeared quasi-distributions;
  \item the \textbf{large-momentum line} (Ji's LaMET and Radyushkin's
        pseudo-PDFs): boosted equal-time correlators matched perturbatively
        onto light-cone distributions.
\end{itemize}
Around them orbit the enabling technologies (Chapter~2), the hard
renormalization problems that made first-principles PDFs possible at all
(Chapter~5), and the two interfaces through which lattice results reach
the rest of physics: global PDF fits (Chapter~7) and machine-learning
samplers rooted in stochastic quantization (Chapter~8).

\section*{How to read the entries}

Each entry opens with a bibliographic card, then answers three questions:
\emph{why} it was selected (with its hub status inside this corpus),
\emph{what} it actually says (physics, with key formulas), and where its
\emph{echoes} appear in the library. Entries are self-contained; formulas
are reproduced in modern notation for orientation, not as facsimiles.

\bigskip
\noindent\emph{A note on honesty.} Citation counts below are counts within
this 99-article corpus (``$\sim$N/99''), rounded conservatively.
Bibliographic data were transcribed from the reference lists of the
library articles themselves and cross-checked against standard records;
any field that could not be confirmed is marked [?].
